% Figure: a 3-coloring of a wheel-ish graph with color classes marked.
% Source for Figure~\ref{fig:coloring}. TikZ only; regenerable by compiling main.tex.
\begin{tikzpicture}[
  c1/.style={circle,draw,fill=accent!55,minimum size=6mm},
  c2/.style={circle,draw,fill=olive!55,minimum size=6mm},
  c3/.style={circle,draw,fill=rulegray!60,minimum size=6mm},
  lab/.style={font=\footnotesize\sffamily,text=slate,align=center}
]
% wheel W5: center + rim of 5
\node[c1] (m) at (0,0) {};
\foreach \i/\col in {0/c2,1/c3,2/c2,3/c3,4/c2} {
  \node[\col] (v\i) at (90+72*\i:1.55cm) {};
}
\foreach \i in {0,...,4} {
  \draw (m) -- (v\i);
  \pgfmathtruncatemacro{\j}{mod(\i+1,5)}
  \draw (v\i) -- (v\j);
}
\node[lab,align=center]{$\chi = 3$: the rim alternates two colours,\\
and the hub takes a third --- three suffice,\\
and two fail since the rim has odd length.};
\node[c1,label=right:{class 1}] at (4.0,0.4) {};
\node[c2,label=right:{class 2}] at (4.0,-0.3) {};
\node[c3,label=right:{class 3}] at (4.0,-1.0) {};
\end{tikzpicture}
